\section{Introduction}
\label{sec:introduction}

\subsection{The Algorithm Selection Problem}

The academic redistricting ensemble literature~\citep{deford2021,mccartan2020,
cannon2022,janson2022} has produced a rich family of algorithms, each
designed with a specific theoretical goal in mind.
DeFord et al.~\citep{deford2021} introduced \recom{} as a general-purpose
Markov chain over the space of valid $k$-district plans.
McCartan and Imai~\citep{mccartan2020} developed Sequential Monte Carlo
(\smc{}) to sample plans with calibrated importance weights, enabling
exact distributional guarantees.
Autry et al.~\citep{autry2021} introduced multi-scale MCMC to address slow
mixing for large $k$.
Each subsequent algorithm in our G-series papers (G.6--G.13) added
further specialisation: compactness optimisation, reversibility, VRA
compliance, or boundary-sensitivity analysis.

For a practitioner---a special master drafting plans for a court, a state
legislative staffer verifying fairness, or an academic studying plan-space
geometry---this proliferation creates a decision problem.
Eight algorithms are implemented; eight sets of CLI flags exist; and each
algorithm has different strengths.
Which one should I run?

\subsection{What This Paper Does}

This paper provides a systematic empirical comparison of all eight
implemented ensemble and search algorithms:

\begin{enumerate}
  \item \textbf{Short-Burst} (\sburst{}, G.6): short \recom{} bursts
        targeting low edge-cut plans~\citep{dellaLibera2026shortburst}.
  \item \textbf{SMC} (\smc{}, G.7): sequential Monte Carlo with
        calibrated importance weights~\citep{dellaLibera2026smc,mccartan2020}.
  \item \textbf{Flip} (\flip{}, G.8): boundary-tract assignment flips for
        local sensitivity analysis~\citep{dellaLibera2026flip}.
  \item \textbf{Forest ReCom} (\fr{}, G.9): reversible recombination
        via spanning forests with MH
        correction~\citep{dellaLibera2026forestRecom}.
  \item \textbf{Merge-Split} (\msp{}, G.10): merge two districts, split
        via random spanning tree cut, approximately
        uniform~\citep{dellaLibera2026mergeSplit}.
  \item \textbf{Multi-scale MCMC} (\ms{}, G.11): interleaved coarse and
        fine chains for large-$k$ mixing~\citep{dellaLibera2026multiscale}.
  \item \textbf{ShortBurstForest} (\sburstf{}, G.12): Short-Burst
        optimisation running on the Forest ReCom chain for compactness
        with distributional
        awareness~\citep{dellaLibera2026shortburstChains}.
  \item \textbf{VRA-Aware ReCom} (\vra{}, G.13): \recom{} with adaptive
        boosts to maintain minority-majority
        districts~\citep{dellaLibera2026vraAware}.
\end{enumerate}

We compare all eight on three states that span the realistic range of
congressional district counts: NC ($k = 14$), WI ($k = 8$), TX ($k = 38$).
For each algorithm--state pair we record four metrics: edge-cut score
(\ec{}), mean Polsby-Popper (\pp{}), partisan seat stability across five
independent runs, and wall-clock runtime.
For chain-based methods we additionally measure lag-100 autocorrelation
of the \ec{} statistic as a proxy for mixing quality.

\subsection{Main Findings}

Our comparison yields five practical conclusions.

\begin{enumerate}
  \item \textbf{No algorithm dominates.}
        \sburst{} achieves the lowest \ec{} (most compact by that metric)
        but concentrates samples near low-\ec{} plans, sacrificing
        distributional correctness.
        \smc{} achieves exact calibration but does not optimise
        compactness.
        \ms{} provides the largest mixing improvement for TX but adds
        overhead for small $k$.

  \item \textbf{Compactness and distribution are in tension.}
        Algorithms that optimise for \ec{} (Short-Burst, ShortBurstForest)
        sample from a concentrated distribution near compact plans, not the
        neutral uniform distribution.
        Practitioners who need a neutral ensemble for litigation should use
        \smc{} or \fr{}, not \sburst{}.

  \item \textbf{Scale matters.}
        For TX ($k = 38$), \ms{} reduces lag-100 autocorrelation by
        approximately 27\% relative to standard \recom{} at comparable
        wall-clock time (a significant but context-specific improvement;
        G.11 reported a 43\% reduction under more favourable initialisation).
        For WI ($k = 8$), the single-scale chain mixes adequately and
        \ms{} adds unnecessary overhead.

  \item \textbf{VRA compliance is orthogonal to optimisation mode.}
        \vra{} achieves distributional properties similar to \fr{}
        when the VRA constraint is not binding, but is the only algorithm
        that reliably maintains minority-majority districts when the
        constraint is active.

  \item \textbf{Decision tree reduces to four questions.}
        A practitioner can select the appropriate algorithm by answering
        four binary questions: (i) Is VRA compliance required? (ii) Is
        $k \geq 20$? (iii) Is a calibrated target distribution required?
        (iv) Is pure compactness optimisation the goal?
        Section~\ref{sec:decision} gives the full tree.
\end{enumerate}

\subsection{Relation to Prior Work}

Cannon et al.~\citep{cannon2022} compare ReCom variants theoretically
but do not include SMC or multi-scale methods.
Janson~\citep{janson2022} analyses the mixing time of random graph
partitioning Markov chains but does not address VRA constraints or
burst optimisation.
Our contribution is purely empirical: we compare all eight algorithms
implemented in \texttt{redist-ensemble} under the same experimental
conditions and translate the results into a practitioner decision framework.

\subsection{Paper Structure}

Section~\ref{sec:setup} describes the experimental setup: three states,
hardware, metrics, and single-run caveats.
Section~\ref{sec:compactness} compares compactness metrics (\ec{} and \pp{})
across all algorithms and states.
Section~\ref{sec:distributional} addresses distributional properties and
mixing speed, including autocorrelation and partisan seat stability.
Section~\ref{sec:decision} presents the decision tree and summary matrix.
Section~\ref{sec:conclusion} concludes with implementation references and
future directions.
